\chapter{\foldertotitle}\label{chp:mathematical_background}
The core of quantum mechanics lies in its underlying mathematical structure. Without a solid grasp of linear algebra we cannot hope to computationally model the behavior of complex quantum systems. While this section does not contain original material, it serves as the crucial foundation for the computational strategies developed and evaluated later in this thesis. We will review the basics of vector spaces, inner products, and linear operators , building up to the concept of Hilbert spaces and Dirac notation, which form the strict mathematical framework for quantum states. Most importantly for our study of nonadiabatic dynamics, we will establish the spectral theory necessary to understand time evolution and rigorously define the adiabatic and diabatic representations. Grasping these specific mathematical definitions is a direct prerequisite for understanding the numerical algorithms implemented in the \textsc{Zinq} code, which is being developed as part of this thesis.\supercite{10.5281/zenodo.18386143}
\inputfolder{sections}
